\documentclass[12pt, a4paper]{article}

\usepackage[utf8]{inputenc}
\usepackage[T1]{fontenc}
\usepackage[brazil]{babel}
\usepackage[margin=2.5cm]{geometry}
\usepackage{parskip}
\usepackage{enumitem}
\usepackage{titlesec}
\usepackage{xcolor}
\usepackage{mdframed}
\usepackage{booktabs}
\usepackage{microtype}
\usepackage{lmodern}

% ---------- Cores e ambientes ----------
\definecolor{destaque}{RGB}{0,80,155}
\definecolor{fundo}{RGB}{240,246,255}
\definecolor{cinzafundo}{RGB}{245,245,245}
\definecolor{bordacinza}{RGB}{180,180,180}

\mdfdefinestyle{sintese}{
    backgroundcolor=fundo,
    linecolor=destaque,
    linewidth=1.2pt,
    leftmargin=0pt, rightmargin=0pt,
    innerleftmargin=10pt, innerrightmargin=10pt,
    innertopmargin=8pt, innerbottommargin=8pt,
    roundcorner=3pt
}
\mdfdefinestyle{tempo}{
    backgroundcolor=cinzafundo,
    linecolor=bordacinza,
    linewidth=0.8pt,
    leftmargin=0pt, rightmargin=0pt,
    innerleftmargin=8pt, innerrightmargin=8pt,
    innertopmargin=5pt, innerbottommargin=5pt,
    roundcorner=2pt
}

% ---------- Seções ----------
\titleformat{\section}{\large\bfseries\color{destaque}}{}{0em}{}[\titlerule]
\titleformat{\subsection}{\normalsize\bfseries}{}{0em}{\thesubsection\quad}

\setlist[itemize]{itemsep=3pt, topsep=4pt}

% ---------- Cabeçalho ----------
\title{\textbf{Roteiro de Apresentação}\\[4pt]
       \large Morten \& Oliveira (2023) ---
       \textit{The Effects of Roads on Trade and Migration}}
\author{Seminário de Doutorado}
\date{Duração estimada: 30 minutos}

\begin{document}
\maketitle
\thispagestyle{empty}

\begin{mdframed}[style=tempo]
\textbf{Distribuição de tempo sugerida:}
Introdução (4 min) $\cdot$
Estratégia Empírica (11 min) $\cdot$
Modelo Estrutural (9 min) $\cdot$
Resultados (4 min) $\cdot$
Conclusão (2 min) $\cdot$
Perguntas buffer (~2 min margem).
\end{mdframed}

\vspace{0.5cm}

% ============================================================
\section{Introdução \hfill \normalfont\small\textit{$\approx$ 4 minutos $\cdot$ slides 1--4}}
% ============================================================

\subsection*{Slide 1 --- Título}

Apresente-se brevemente e anuncia o artigo. Não gaste tempo aqui.

\begin{mdframed}[style=sintese]
``O artigo que vou apresentar é de Melissa Morten e Joana Oliveira, publicado em 2023.
O tema central é o papel das estradas sobre comércio e migração no Brasil, usando a
construção de Brasília como experimento natural.''
\end{mdframed}

\subsection*{Slide 2 --- Motivação}

\begin{itemize}
    \item Destaque que a literatura sobre infraestrutura concentrou atenção no comércio de bens, deixando de lado o canal de migração.
    \item As duas perguntas do artigo são complementares: a primeira é positiva (estradas promovem migração?), a segunda é normativa (quanto isso importa para o bem-estar?).
\end{itemize}

\begin{mdframed}[style=sintese]
``A pergunta central do artigo é dupla. Primeiro, uma questão positiva: estradas facilitam a
migração interna, além do comércio? Segundo, uma questão normativa: se facilita os dois,
qual canal gera mais bem-estar? A contribuição do artigo é responder as duas ao mesmo tempo,
com um único arcabouço empírico e estrutural.''
\end{mdframed}

\subsection*{Slide 3 --- Contribuições}

\begin{itemize}
    \item Contribuição 1 é essencialmente identificação causal --- o instrumento MST é o ponto forte metodológico.
    \item Contribuição 2 é o modelo estrutural --- ele permite ir além das elasticidades e falar em bem-estar.
\end{itemize}

\begin{mdframed}[style=sintese]
``O artigo tem duas contribuições distintas que se complementam. A parte empírica estabelece que estradas
promovem migração de forma causal, usando um instrumento que vou detalhar daqui a pouco. A parte estrutural
usa essas estimativas para decompor os ganhos de bem-estar entre o canal de comércio e o canal de migração.''
\end{mdframed}

\subsection*{Slide 4 --- Sumário executivo}

Apresente os números principais sem entrar em detalhes. Eles voltarão mais tarde com contexto.

\begin{mdframed}[style=sintese]
``Para adiantar o resultado: a rede rodoviária de Brasília aumentou o bem-estar médio em 2,8\%.
Três quartos desse ganho vêm do canal comercial, um quarto do canal migratório. E, surpreendentemente,
a localização de Brasília foi ela mesma vantajosa: uma rede alternativa centrada no Rio teria gerado
0,2 a 0,3 ponto percentual a menos de bem-estar.''
\end{mdframed}

\vspace{0.3cm}

% ============================================================
\section{Estratégia Empírica \hfill \normalfont\small\textit{$\approx$ 11 minutos $\cdot$ slides 5--13}}
% ============================================================

\subsection*{Slide 6 --- Dados}

\begin{itemize}
    \item Dados de migração vêm dos Censos: cada observação é um par (origem, destino, ano).
    \item Dados de comércio vêm de Anuários Estatísticos, cobrem mais anos que os censos.
    \item A malha rodoviária é reconstruída historicamente com SIG e algoritmo de \textit{fast marching} para tempo de viagem.
    \item Ressalte que tempo de viagem (não distância) é a variável de fricção. Isso é mais preciso porque a qualidade da estrada varia.
\end{itemize}

\begin{mdframed}[style=sintese]
``Uma decisão metodológica importante é usar \textit{tempo de viagem} entre centroides estaduais
em vez de distância geográfica. Isso captura o fato de que a construção e pavimentação de rodovias
não diminui a distância, mas diminui substancialmente o tempo de deslocamento.''
\end{mdframed}

\subsection*{Slide 7 --- Modelo de gravidade (eq.\ 1)}

\begin{itemize}
    \item A equação de gravidade é o padrão da literatura de comércio internacional, aqui aplicada a fluxos internos.
    \item Os efeitos fixos de par $\gamma_{od}$ absorvem tudo que não varia no tempo: distância, fronteiras, cultura. O que identifica $\beta$ são as \textit{mudanças} no tempo de viagem.
    \item $\beta < 0$ esperado: mais tempo de viagem $\Rightarrow$ menos migração/comércio.
\end{itemize}

\vspace{2cm}

\begin{mdframed}[style=sintese]
``O modelo de gravidade aqui é saturado de efeitos fixos. O par de efeitos fixos de origem-ano e destino-ano
absorve todos os choques de cada estado em cada período. O que sobra para identificar $\beta$ é a variação
no tempo de viagem \textit{dentro} de cada par $(o,d)$ ao longo do tempo. Isso é uma especificação bem exigente.''
\end{mdframed}

\subsection*{Slide 8 --- Identificação: o problema}

\begin{itemize}
    \item O problema de endogeneidade aqui é clássico: o governo pode construir estradas \textit{em resposta} a condições econômicas.
    \item O viés é positivo: choques negativos geram tanto menos migração/comércio quanto mais investimento em infraestrutura. O MQO captura isso e subestima o efeito real das estradas.
\end{itemize}

\begin{mdframed}[style=sintese]
``O problema é que a alocação de estradas não é aleatória. Se o governo tende a construir estradas
onde a economia está em dificuldade, a correlação entre estradas e comércio/migração vai ser menor
do que o efeito real. O MQO vai \textit{subestimar} o impacto causal --- e é exatamente isso que vamos ver nos resultados.''
\end{mdframed}

\subsection*{Slides 9--10 --- O instrumento MST}

\begin{itemize}
    \item Explique com calma os quatro passos: centrar em Brasília, dividir em oito fatias, menor caminho em cada fatia, unir.
    \item A intuição chave: a rede predita pelo MST é determinada por geometria e pela localização da capital --- não por demanda econômica entre pares de estados.
    \item O mapa (Fig.\ 1) ajuda muito: mostre que a rede real e a MST se sobrepõem parcialmente, mas a MST corta por caminhos que não seriam escolhidos por demanda econômica pura.
\end{itemize}

\begin{mdframed}[style=sintese]
``A ideia do instrumento é engenhosa. A pergunta que o MST responde é: \textit{se o único objetivo
fosse conectar Brasília às capitais estaduais pelo caminho mais curto em cada direção cardinal,
como seria a rede?} Essa rede hipotética gera variação no tempo de viagem entre pares de estados
que é determinada pela geometria do país e pela localização da capital --- não por qualquer
demanda econômica bilateral. É essa variação que identifica o efeito causal das estradas.''
\end{mdframed}

\subsection*{Slide 11 --- Especificação com instrumento (eq.\ 2)}

\begin{itemize}
    \item O instrumento é invariante no tempo porque a rede MST é uma característica geográfica fixa de cada par.
    \item Os $\alpha_t$ variam por ano: permitem testar em \textit{que momento} os efeitos aparecem.
    \item Normalização em 1950 (pré-Brasília) permite teste de pré-tendências.
\end{itemize}

\vspace{1cm}

\begin{mdframed}[style=sintese]
``O instrumento em si não varia no tempo --- ele é uma propriedade geométrica fixa do par $(o,d)$.
O que varia é \textit{quando} esse instrumento se torna relevante: antes de Brasília existir, a rede MST
não tinha correspondência real. Depois de 1960, a rede efetiva começa a se aproximar da rede predita.
Os $\alpha_t$ capturam exatamente essa dinâmica.''
\end{mdframed}

\subsection*{Slide 12 --- Efeitos ao longo do tempo (Fig.\ 2)}

\begin{itemize}
    \item Três painéis: tempo de viagem, migração, comércio. Todos normalizados em 1950.
    \item Coeficiente de 1940 não significativo: sem pré-tendência. Isso valida o instrumento.
    \item Após 1960: queda brusca no tempo de viagem, alta em migração e comércio.
    \item Coeficientes de 1960: +0,4 migração, +1,4 comércio para uma redução de 10 p.p. MST-predita.
\end{itemize}

\begin{mdframed}[style=sintese]
``O teste de pré-tendências é essencial para qualquer instrumento desse tipo. Se em 1940 o coeficiente
já fosse significativo, indicaria que o MST estava capturando algo relacionado a tendências econômicas
pré-existentes. O fato de o coeficiente de 1940 não diferir de zero é o sinal verde para a validade do instrumento.
A quebra estrutural em 1960 é exatamente o que esperamos: Brasília foi inaugurada em abril de 1960.''
\end{mdframed}

\subsection*{Slide 13 --- Tabela 1 e principais achados}

\begin{itemize}
    \item Compare MQO vs.\ VI: o VI sistematicamente maior em módulo confirma subestimação do MQO.
    \item Colunas PPML com função de controle: robusto a zeros --- importante porque vários pares têm fluxo zero.
    \item Destacar robustez: excluir pares com Goiás não altera os resultados.
\end{itemize}

\begin{mdframed}[style=sintese]
``O resultado mais importante da tabela é a diferença entre MQO e VI. Para migração, o MQO dá $-0{,}33$
e o VI dá $-1{,}75$ --- cinco vezes maior em módulo. Para comércio, $-0{,}78$ vs.\ $-1{,}90$.
Isso confirma o viés previsto: o governo realmente construiu mais estradas onde a economia sofria,
comprimindo artificialmente a correlação observada. O efeito real das estradas é substancialmente maior.''
\end{mdframed}

\vspace{0.3cm}

% ============================================================
\section{Modelo Estrutural \hfill \normalfont\small\textit{$\approx$ 9 minutos $\cdot$ slides 14--21}}
% ============================================================

\subsection*{Slide 15 --- Estrutura do modelo}

\begin{itemize}
    \item Dois blocos emprestados da literatura: Eaton-Kortum (2002) para comércio, Monte et al.\ (2018) para migração.
    \item O que é novo é a combinação dos dois em equilíbrio geral, permitindo a decomposição.
    \item Destaque os dois parâmetros de fricção: $\tau_{odt}$ (custo iceberg de comércio) e $\kappa_{odt}$ (custo de migração). Ambos dependem do tempo de viagem --- esse é o elo entre a evidência empírica e o modelo.
\end{itemize}

\begin{mdframed}[style=sintese]
``O modelo combina dois ingredientes conhecidos: a estrutura Fréchet de Eaton-Kortum para comércio
e a decisão de migração discreta com custos de Monte et al. O que os autores fazem de novo é
colocar os dois num mesmo equilíbrio geral e disciplinar os parâmetros de fricção com as estimativas
do instrumental da parte empírica.''
\end{mdframed}

\subsection*{Slide 16 --- Utilidade do consumidor (eq.\ 3)}

\begin{itemize}
    \item A equação de utilidade indireta é o ponto de partida do bem-estar.
    \item Estradas afetam $V_{dt}$ por dois caminhos simultâneos: reduzem $P_{dt}$ (via comércio mais barato) e reduzem $\kappa_{odt}$ (via migração mais fácil). Esses são os dois canais da decomposição.
\end{itemize}

\begin{mdframed}[style=sintese]
``A utilidade indireta captura tudo que importa para o bem-estar de quem vive em $d$:
salários, nível de preços de bens tradable e custo de moradia. As estradas entram por $P_{dt}$
--- bens importados ficam mais baratos --- e por $\kappa_{odt}$ --- é mais fácil migrar para
um destino melhor. Os dois canais operam simultaneamente, por isso precisamos do modelo para separá-los.''
\end{mdframed}

\subsection*{Slide 17 --- Decisão de migração (eq.\ 4)}

\begin{itemize}
    \item Estrutura logit Fréchet: a probabilidade de migrar para $d$ aumenta com $V_{dt}$ e cai com $\kappa_{odt}$.
    \item $\epsilon$ é a elasticidade de migração. Estimado em 4,5 --- heterogeneidade moderada nas preferências.
    \item A fórmula fechada é análoga à de Eaton-Kortum: a divisão pelo denominador $\Phi_{ot}$ garante que as probabilidades somem 1.
\end{itemize}

\subsection*{Slide 18 --- Comércio e preços (eq.\ 5)}

\begin{itemize}
    \item Estrutura idêntica à de migração, com os parâmetros trocados: $\theta$ em vez de $\epsilon$, $\tau$ em vez de $\kappa$.
    \item A analogia entre as duas equações é didaticamente importante: o mesmo framework cobre bens e pessoas.
    \item $P_{dt}$ cai com $\Theta_{dt}$: quanto mais fontes baratas/produtivas $d$ acessa, menor seu nível de preços.
\end{itemize}

\begin{mdframed}[style=sintese]
``Uma elegância do modelo é que as equações de comércio e migração têm a mesma estrutura matemática.
Em ambas, a fração de bens (ou de pessoas) que flui de $o$ para $d$ é proporcional à atratividade
de $d$ dividida pela soma de atratividade de todos os destinos. A diferença está nos parâmetros e
nas fricções: $\tau$ para bens, $\kappa$ para pessoas.''
\end{mdframed}

\subsection*{Slide 19 --- Equilíbrio de salários e aluguéis (eqs.\ 6--7)}

\begin{itemize}
    \item Equação de salários é market-clearing do mercado de bens: renda de $o$ deve igualar exportações totais de $o$.
    \item Equação de aluguéis: aluguel cresce com população e salários --- aglomeração pressiona moradia.
    \item Destaque a tensão: estradas aumentam salários e população, mas também sobem aluguéis. O bem-estar líquido é ambíguo \textit{a priori} --- é por isso que precisamos do modelo.
\end{itemize}

\begin{mdframed}[style=sintese]
``O mercado de habitação é importante para disciplinar os ganhos de bem-estar. Se a oferta de
moradia for inelástica --- e $1/\epsilon_r = 1{,}7$ sugere que é --- aluguéis vão subir com
o influxo de trabalhadores e absorver parte do ganho salarial. Isso cria um mecanismo de
compressão dos ganhos que é heterogêneo: estados que atraem muita migração sofrem mais
pressão de aluguéis.''
\end{mdframed}

\subsection*{Slide 20 --- Efeitos de bem-estar (eqs.\ 8--10)}

\begin{itemize}
    \item O método ``hat'' ($\widehat{X} = X_{\mathrm{novo}} / X_{\mathrm{ant}}$) é padrão em comércio internacional desde Dekle et al.\ (2008). Permite calcular mudanças de equilíbrio sem precisar do nível absoluto dos dados.
    \item Bem-estar agregado é média ponderada das variações de $\Phi_{ot}^{1/\epsilon}$ sobre origens.
    \item A decomposição sai naturalmente: fixar $\hat\kappa = 1$ (sem mudança nos custos de migração) mostra o ganho puramente comercial, e vice-versa.
\end{itemize}

\begin{mdframed}[style=sintese]
``A decomposição é operacionalizada pelo método hat: você resolve o equilíbrio duas vezes.
Na primeira, permite que tanto $\tau$ quanto $\kappa$ mudem. Na segunda, mantém $\kappa$ fixo
em 1 e deixa só $\tau$ variar. A diferença entre os dois cenários é o quanto do ganho total
vem exclusivamente da melhora na mobilidade de pessoas.
Isso não seria possível sem o modelo estrutural --- a parte econométrica sozinha não consegue separar os canais.''
\end{mdframed}

\subsection*{Slide 21 --- Parâmetros estruturais}

\begin{itemize}
    \item $\theta = 4$ é externo (Simonovska e Waugh, 2014) --- padrão na literatura.
    \item $\epsilon = 4{,}5$ e $1/\epsilon_r = 1{,}7$ são estimados pelos próprios autores via instrumento Bartik.
    \item Ressalte que os parâmetros são modestos o suficiente para não gerar resultados mecânicos: o resultado de 76/24 não está embutido nas hipóteses.
\end{itemize}

\vspace{0.3cm}

% ============================================================
\section{Resultados \hfill \normalfont\small\textit{$\approx$ 4 minutos $\cdot$ slides 22--24}}
% ============================================================

\subsection*{Slide 22 --- Decomposição do bem-estar}

\begin{itemize}
    \item Número central: +2,8\% de bem-estar médio com a rede de Brasília vs.\ ausência de estradas.
    \item A decomposição 76/24 é o resultado que justifica o título do artigo: \textit{roads affect both, but trade dominates}.
    \item Por que o canal de comércio domina? Porque $\theta = 4$ amplifica muito as reduções em $\tau$, e os custos de comércio caíram 9\% contra 8\% nos custos de migração --- uma diferença pequena que o parâmetro amplifica.
\end{itemize}

\begin{mdframed}[style=sintese]
``O canal de comércio domina por duas razões. Primeira, a elasticidade de comércio $\theta = 4$
é maior que a elasticidade de migração $\epsilon = 4{,}5$... espera, os parâmetros são próximos.
O que diferencia é a \textit{magnitude} dos canais na utilidade: o índice de preços afeta diretamente
o bem-estar de toda a população do destino, enquanto o ganho de migração beneficia diretamente
apenas quem de fato migra. Como a maioria das pessoas não muda de estado, o canal comercial
tem uma base muito maior.''
\end{mdframed}

\subsection*{Slide 23 --- Heterogeneidade espacial}

\begin{itemize}
    \item A amplitude 1\%--15\% mostra que os ganhos são muito desiguais entre estados.
    \item A comparação com mobilidade perfeita é o diagnóstico: se todos migrassem livremente, os ganhos seriam equalizados. A heterogeneidade observada é inteiramente gerada pelos custos de migração.
    \item O contrafactual Rio de Janeiro: ressalte que a diferença de 0,2 p.p.\ é pequena em termos absolutos, mas significativa considerando que estamos comparando duas localizações de capital igualmente prováveis.
\end{itemize}

\begin{mdframed}[style=sintese]
``O resultado de heterogeneidade espacial tem implicação de política clara: simplesmente
construir estradas não é suficiente para equalizar bem-estar entre regiões. Os custos
de migração agem como fricção que impede que os trabalhadores de estados menos conectados
se beneficiem plenamente da integração. Políticas complementares de mobilidade --- subsídios
de relocação, acesso a informação sobre mercados de trabalho --- teriam valor expressivo.''
\end{mdframed}

\subsection*{Slide 24 --- Resultados principais (resumo da Tabela 1)}

Caso o tempo ainda permita, você pode revisitar rapidamente os achados econométricos antes de concluir. Caso contrário, siga direto para a conclusão.

\vspace{0.3cm}

% ============================================================
\section{Conclusão \hfill \normalfont\small\textit{$\approx$ 2 minutos $\cdot$ slide 25}}
% ============================================================

\subsection*{Slide 25 --- Conclusão}

\begin{itemize}
    \item Recapitule as duas contribuições: evidência causal de que estradas promovem migração + decomposição do bem-estar.
    \item Implicação para a literatura: a desconsideração do canal de migração em estudos anteriores de infraestrutura provavelmente subestimou os ganhos totais de bem-estar.
    \item Deixe aberto para perguntas; o artigo tem várias extensões (heterogeneidade de trabalhadores por qualificação, efeitos de longo prazo) que podem surgir na discussão.
\end{itemize}

\begin{mdframed}[style=sintese]
``O artigo deixa uma lição metodológica importante: quando estimamos o impacto de infraestrutura
usando apenas dados de comércio, capturamos no máximo 76\% do efeito de bem-estar. O canal
migratório, embora menor, não é negligenciável --- e é ele que explica por que os ganhos não
são uniformes no espaço. Para políticas de desenvolvimento regional, ignorar esse canal significa
subestimar tanto os benefícios da infraestrutura quanto sua capacidade de ampliar desigualdades regionais.''
\end{mdframed}

\vspace{0.5cm}
\hrule
\vspace{0.3cm}

\subsection*{Pontos que costumam gerar perguntas}

\begin{itemize}
    \item \textbf{Validade externa do instrumento:} o MST é específico ao Brasil e à decisão de construir Brasília. Como generalizar? O artigo argumenta que o mecanismo --- fricção de custo reduzida por estradas --- é geral; o instrumento é o meio de identificá-lo.
    \item \textbf{Por que o MQO subestima (e não superestima)?} Porque o governo construiu estradas \textit{onde a economia estava mal}, criando correlação espúria negativa entre infraestrutura e atividade econômica.
    \item \textbf{O papel de Goiás:} Como Brasília fica em Goiás, pares envolvendo esse estado poderiam capturar efeitos da capital em si. A robustez de excluir Goiás descarta isso.
    \item \textbf{Habitação inelástica:} Um crítico pode perguntar se $\epsilon_r = 1/1{,}7$ é razoável para o Brasil dos anos 1960--2000. Resposta: é estimado pelos dados, não calibrado ad hoc.
    \item \textbf{Migração vs.\ comércio: complementares ou substitutos?} No modelo, são complementares: estradas que facilitam comércio também atraem trabalhadores, e mais trabalhadores geram mais produção e mais comércio.
\end{itemize}

\end{document}
